% AUTORIDAD DE FUENTE:
% expose_VIII_body.tex, líneas 1917--1974; índices de escaneo 295--298;
% folios de fuente 284--287. Las correcciones materiales de la capa
% francesa se conservan reversiblemente mediante \sisi.

\medskip
\noindent 5.7. Sean $S$ un esquema, $T$ un subesquema cerrado y $Z$
un conjunto. Se llama \underline{esquema constante sobre $S$, truncado
en $T$, de valor $Z$}, y se denota por $Z_{S,T}$, al esquema que se
obtiene tomando la suma de copias $S_z$ de $S$, una por cada $z\in Z$,
y pegando todas estas copias a lo largo del subesquema abierto
$U=S\setminus T$. Así,
\begin{equation}\tag{5.7.1}
 Z_{S,T}\times_S U\simeq U,
 \qquad
 Z_{S,T}\times_S T\simeq Z_T,
\end{equation}
y $Z_{S,T}$ es étale sobre $S$. Para un $S$-esquema variable $S'$
existe un isomorfismo, funtorial en $S'$,
\begin{equation}\tag{5.7.2}
 \Hom_S(S',Z_{S,T})\simeq\Hom_T(S'_T,Z_T).
\end{equation}
La aplicación es la restricción a $S'_T$, usando el segundo
isomorfismo de (5.7.1); dejamos al lector comprobar que es biyectiva.
Si

%%%%%%%%%%  scan idx 296  |  folio 285  %%%%%%%%%%

\[
 i:T\longrightarrow S
\]
denota la inclusión, el isomorfismo funtorial (5.7.2) también puede
escribirse como un isomorfismo de haces
\begin{equation}\tag{5.7.3}
 Z_{S,T}\simeq i_*(Z_T).
\end{equation}
Esta fórmula, o directamente la definición, muestra que, al variar
$Z$, $Z_{S,T}$ depende funtorialmente de $Z$ y que este funtor conmuta
con los límites proyectivos finitos. En particular, transforma grupos
en esquemas en grupos y grupos conmutativos en esquemas en grupos
conmutativos, como también resulta claro de la forma
\sisi{(2.7.2)}{(5.7.2)} del funtor representado por $Z_{S,T}$.

También es inmediato a partir de la definición, o de (5.7.3), que la
formación de $Z_{S,T}$ conmuta con todo cambio de base $f:S'\to S$.
Esto da un isomorfismo, funtorial en $Z$,
\begin{equation}\tag{5.7.4}
 Z_{S,T}\times_S S'\simeq Z_{S',T'},
 \qquad T'=T\times_S S'=T_{S'}.
\end{equation}

\medskip
\noindent 5.8. Estudiaremos los torsors $E$ bajo un esquema en grupos
de la forma $Z_{S,T}$, donde $Z$ es un grupo, que supondremos
conmutativo para simplificar. Existe un funtor natural
\begin{equation}\tag{5.8.0}
 E\longmapsto E_T
\end{equation}
de la categoría de torsors bajo $Z_{S,T}$ a la categoría de torsors
bajo $Z_T$. Afirmamos que este funtor es una
\underline{equivalencia de categorías}. Su plena fidelidad significa
que, para dos torsors $E,F$ bajo $\sisi{G_{S,T}}{Z_{S,T}}$, la
aplicación
\[
 \Hom(E,F)\longrightarrow\Hom(E_T,F_T)
\]
es biyectiva. Por descenso fpqc, esto se reduce inmediatamente al caso
en que $E$ y $F$ son triviales,

%%%%%%%%%%  scan idx 297  |  folio 286  %%%%%%%%%%

y, por tanto, a $E=F=\sisi{G_{S,T}}{Z_{S,T}}$, donde la afirmación es
precisamente la biyectividad de (5.7.2) para $S'=S$.

Queda demostrar que (5.8.0) es esencialmente sobreyectivo. Sea $E_T$
un torsor bajo $Z_T$. Ya observamos en la demostración de 5.1 que
$E_T$ es representable y localmente trivial para la topología étale.
Basta, pues, demostrar que
\[
 H^1(S_{\mathrm{\acute{e}t}},Z_{S,T})
 \longrightarrow H^1(T_{\mathrm{\acute{e}t}},Z_T)
\]
es biyectiva. Esto resulta de (5.7.3) y de la relación
\[
 R^1i_{\mathrm{\acute{e}t},*}(Z_T)=0,
\]
que es el caso particular SGA 4 VIII 5.6.

\medskip
\noindent 5.8.1. Sea ahora $P$ un esquema en grupos conmutativo sobre
$S$. Queremos determinar la categoría de extensiones, como haces de
grupos conmutativos, de $P$ por $Z_{S,T}$. Para ello utilizamos la
descripción de VII 1.1.6 y 1.3 en términos de torsors sobre $P$,
$P\times P$ y $P\times P\times P$, bajo las imágenes inversas de
$Z_{S,T}$. Por (5.7.4), estas imágenes inversas vuelven a ser de la
misma forma que $Z_{S,T}$. Aplicando a esos torsors el resultado
precedente, se ve que corresponden a torsors sobre $P_T$,
$P_T\times_TP_T$ y $P_T\times_TP_T\times_TP_T$, bajo las imágenes
inversas de $Z_T$. Obtenemos, por tanto, lo siguiente.

\medskip
\noindent\underline{Proposición 5.9.} \underline{Sean $S$ un esquema,
$T$ un subesquema cerrado, $\sisi{G}{Z}$ un grupo conmutativo,
$\sisi{G_{S,T}}{Z_{S,T}}$ el correspondiente esquema en grupos
constante truncado sobre $S$, definido en 5.7, y $P$ un esquema en
grupos conmutativo sobre $S$. Entonces la restricción a $T$ induce una
equivalencia entre la categoría}

%%%%%%%%%%  scan idx 298  |  folio 287  %%%%%%%%%%

\underline{de haces de grupos conmutativos que son extensiones de $P$
por $Z_{S,T}$ y la categoría de haces de grupos conmutativos que son
extensiones de $P_T$ por $Z_T$.}

Una demostración esencialmente idéntica da el resultado siguiente.

\medskip
\noindent\underline{Corolario 5.10.} \underline{Con la notación de
5.9, sea $Q$ otro esquema en grupos conmutativo sobre $S$. Entonces la
restricción a $T$ induce una equivalencia entre la categoría de
biextensiones de $(P,Q)$ por $Z_{S,T}$ y la categoría de
biextensiones de $(P_T,Q_T)$ por $Z_T$.}
